% ===================================================================
\chapter{Pruebas, calidad y documentación}

\section{Pruebas automáticas}
\begin{terminal}
uv run pytest -m "not slow"          # pruebas rápidas (~10 s)
uv run pytest                        # todas, incluidas las lentas (~1 min)
uv run pytest tests/test_fase3.py    # solo una fase
uv run pytest -k kuramoto            # solo las que contienen "kuramoto" en el nombre
\end{terminal}
Muchas pruebas comparan el código con resultados analíticos conocidos: la solución
exacta del sistema lineal, la conservación bajo difusión, el valor de Kuramoto en la red
completa, Eisenberg–Noe resuelto a mano. \texttt{tests/test\_app.py} comprueba que cada
página de la aplicación se ejecute sin errores.

\section{Estilo del código}
\begin{terminal}
uv run ruff check src app tests          # errores y estilo
uv run ruff format src app tests         # formato automático
\end{terminal}

\section{Integración continua}
Al subir cambios a GitHub, el flujo \ruta{.github/workflows/ci.yml} instala las
dependencias, corre ruff y pytest y construye la imagen de Docker. Un cambio que rompa
una prueba aparece marcado en rojo en el repositorio.

\section{Reproducibilidad}
\begin{itemize}
\item Todos los parámetros están en \ruta{configs/}.
\item Todas las semillas son explícitas: la misma configuración produce los mismos
resultados, bit a bit, con cualquier número de procesos.
\item Las tablas se validan contra un esquema de columnas, tipos y rangos antes de
guardarse o de entrar a un modelo.
\item \texttt{uv.lock} fija la versión exacta de cada dependencia.
\end{itemize}

\section{Documentos y cómo regenerarlos}
\begin{longtable}{L{0.4\textwidth}L{0.52\textwidth}}
\toprule
\textbf{Documento} & \textbf{Contenido} \\
\midrule
\endhead
\ruta{v1/pulso\_guia\_completa.pdf} & Fundamento completo desde cero. \\
\ruta{v1/pulso\_presentacion.pdf} & Presentación completa (56 diapositivas). \\
\ruta{v2/pulso\_presentacion\_v2.pdf}, guion & Presentación breve y su guion. \\
\ruta{v2/pulso\_manual\_app.pdf} & Manual de la aplicación (incluido en este manual). \\
\ruta{v3/pulso\_presentacion\_v3.pdf} & Presentación de 30 minutos. \\
\ruta{v3/pulso\_guion\_v3.pdf}, \texttt{\_15min} & Guiones de 30 y 15 minutos. \\
\ruta{ESTUDIO.pdf} & Plan de estudio de los temas del proyecto. \\
\ruta{MANUAL\_DE\_USUARIO.pdf} & Este manual. \\
\bottomrule
\end{longtable}
Todas las rutas son relativas a \ruta{docs/Presentaciones/}. Para regenerarlos se
necesita TeX Live con \texttt{lualatex}, \texttt{latexmk} y el tema \emph{metropolis}:
\begin{terminal}
uv run pulso credit configs/credit.yaml              # resultados de la tarea 6
PYTHONPATH=src uv run python docs/latex/figuras.py    # figuras de la tarea 6
bash docs/latex/compilar.sh                           # todos los PDF
\end{terminal}

% ===================================================================
\chapter{Solución de problemas}
\label{cap:problemas}

\begin{longtable}{L{0.38\textwidth}L{0.56\textwidth}}
\toprule
\textbf{Síntoma} & \textbf{Causa y solución} \\
\midrule
\endhead
\texttt{ModuleNotFoundError: No module named 'pulso'} (macOS) & macOS oculta un archivo
de la instalación cuando el proyecto está dentro de \texttt{\textasciitilde/Documents}.
Ejecute \texttt{chflags -R nohidden .venv}, o anteponga \texttt{PYTHONPATH=src} a cada
comando. \\
\texttt{No existe el directorio de configuración} & El comando no se corrió desde la raíz
del proyecto. Cámbiese a esa carpeta o defina \texttt{PULSO\_CONFIG\_DIR}. \\
\texttt{No se encontró el archivo de configuración} & La ruta del YAML está mal escrita.
Use \texttt{configs/<nombre>.yaml}. \\
\emph{Port 8501 is already in use} & Otra aplicación usa el puerto. Agregue
\texttt{--server.port=8502}. \\
\emph{No hay dataset} (Explorar escenarios) & Ejecute \texttt{uv run pulso sweep
configs/sweep.yaml}. \\
\emph{No hay modelos entrenados} (Mapa de riesgo) & Ejecute \texttt{uv run pulso train
configs/train.yaml}. \\
\emph{No hay resultados del experimento} (páginas 8 y 9) & Ejecute \texttt{uv run pulso
credit configs/credit.yaml}. \\
\emph{Red demasiado grande para simulación en vivo} & La red tiene más de 300 entidades.
Redúzcala en \emph{Construir red}. \\
La página no cambia al mover un control & El control está en un formulario: presione su
botón azul y espere a que termine el indicador de ejecución. \\
Los valores volvieron a los iniciales & Se recargó la pestaña del navegador. \\
Cierre inesperado al usar XGBoost y PyTorch juntos & Ambos traen su propia librería
OpenMP. El proyecto ya carga XGBoost primero y limita PyTorch a un hilo; si escribe código
propio, importe \texttt{xgboost} antes que \texttt{torch}. \\
Error de YAML al leer la configuración & Revise la sangría (solo espacios) y los dos
puntos después de cada clave. \\
\bottomrule
\end{longtable}

% ===================================================================
\chapter{Glosario rápido}

\begin{longtable}{L{0.24\textwidth}L{0.7\textwidth}}
\toprule
\textbf{Término} & \textbf{Significado} \\
\midrule
\endhead
$\alpha$ & Velocidad con que una entidad corrige su déficit. \\
$\beta$ & Intensidad del contagio de déficit. \\
$\gamma$ & Déficit que genera el ingreso desfasado. \\
$\Delta$ & Dispersión de los ritmos naturales de pago. \\
$\lambda$ & Déficit que genera perder un pagador. \\
$D_0$ & Intensidad de la redistribución de liquidez. \\
$K$ & Fuerza con que se ajustan los ritmos de pago. \\
$R$ & Índice de amplificación; menor que 1, la red amortigua. \\
$\sigma$, $\tau_\text{rec}$ & Estabilidad y tiempo de recuperación. \\
$r$, $\hat r$ & Sincronía real y estimada desde transacciones. \\
AUC & Capacidad de un modelo para ordenar a quienes caen por encima de quienes no. \\
CLI & Interfaz de línea de comandos: el comando \texttt{pulso}. \\
Docker & Empaqueta el programa para que corra igual en cualquier equipo. \\
GCN & Red neuronal que trabaja sobre grafos. \\
Parquet & Formato de archivo para tablas, comprimido y con tipos fijos. \\
PD & Probabilidad de impago a doce meses. \\
SHAP & Método que explica cuánto aporta cada variable a una predicción. \\
SMMLV & Salario mínimo mensual legal vigente. \\
Streamlit & Biblioteca con que está hecha la aplicación web. \\
uv & Gestor de entornos y dependencias de Python. \\
XGBoost & Modelo de árboles de decisión entrenados en secuencia. \\
YAML & Formato de texto de los archivos de configuración. \\
\bottomrule
\end{longtable}
